\section{Синтез формы занавеса}

Пусть $u \in [0;\, \ell]$ "--- координата вдоль основного направления занавеса, $w(u, t)$ "--- отклонение основания в момент времени $t$. Функция $w$ должна удовлетворять условию $|w| \leqslant A$, иметь характерную длину складки $\lambda$, определяться значением $r_0$ и непрерывно изменяться во времени; функцией того же вида задаётся множитель яркости вдоль занавеса.

\subsection{Гармоническая модель}

Отклонение задаётся суммой $J$ синусоид, как в рядах Фурье~\cite[с.~250]{gonzalez2012}:
\begin{equation}
	w(u, t) = \sum_{j=1}^{J} a_j \sin\left( \frac{2\pi u}{\lambda_j} + \omega_j t + \varphi_j \right), \qquad \sum_{j=1}^{J} a_j = A,
	\label{eq:harmonic}
\end{equation}
где периоды $\lambda_j$ выбираются вблизи $\lambda$, а фазы $\varphi_j$ и частоты $\omega_j$ порождаются генератором псевдослучайных чисел. Такие складки выглядят регулярными, а со временем лишь сдвигаются вдоль занавеса.

\subsection{Шум Перлина}

Шум Перлина "--- псевдослучайная функция со значениями в отрезке $[-1;\, 1]$, заданная в узлах целочисленной решётки и гладко интерполируемая между ними~\cite{hill2002}:
\begin{equation}
	w(u, t) = A \cdot \mathrm{noise}\left( u / \lambda,\ t / \tau \right),
	\label{eq:perlin}
\end{equation}
где $\tau$ "--- постоянный масштаб времени. Параметры $A$ и $\lambda$ входят в формулу непосредственно, картина не повторяется, а форма меняется плавно.

\subsection{Выбор метода синтеза формы занавеса}

Основание вычисляется не более чем в $\ell / \delta \leqslant 256$ точках на занавес, поэтому трудоёмкость обоих методов пренебрежимо мала; сравнение приведено в таблице~\ref{tab:synthesis}.

\begin{table}[H]
	\caption{Сравнение методов синтеза формы занавеса}
	\label{tab:synthesis}
	\begin{tabular}{|L{0.44}|C{0.27}|C{0.27}|}
		\hline
		Критерий & Гармоническая модель & Шум Перлина \\
		\hline
		Параметры сверх задания & $J$ & нет \\
		\hline
		Повторяемость картины & есть & нет \\
		\hline
		Изменение во времени & сдвиг составляющих & изменение формы \\
		\hline
	\end{tabular}
\end{table}

Выбран шум Перлина~\eqref{eq:perlin}: он использует только параметры задания и не даёт повторяющихся складок.